\chapter{Vågutbredning och antenner}

I detta kapitlet tittar vi närmare på olika sorters vågutbredning och
hur de fungerar, fördelar och nackdelar och hur man kan använda dem
för att kommunicera på olika avstånd. Vi tar också upp lite om hur
rymdvädret kan påverka radiokommunikationen samt en del enkel
antennlära.

\vfill

\minitoc

\clearpage


\label{vågutbredning}
\index{vågutbredning}


\section{Vågutbredning}
\index{vågutbredning}

Vågutbredning är något som är väldigt viktigt för radiokommunikation och
beroende på vilken frekvens som används kan vågutbredningen se mycket olika ut
för olika radiosändningar. För VHF och uppåt är det i princip bara markvåg men
för lägre frekvenser kan också rymdvågsutbredning ske.

\subsection{Markvåg}
\index{markvåg}
\index{vågutbredning!markvåg}

Markvågen är den vanligaste radioutbredningen vi tänker oss, här får man hjälp
av att placera antennan över terrängen och på hustak för att få så fri sikt
som möjligt vilket ökar längden på markvågsutbredningen. Det är också den
utbredning som handhållna radioapparater nyttjar men i och med att man som
regel då kommer vara mycket nära marken blir oftast utbredningen begränsad.

I närheten av vatten kan markvågen ofta gå längre än till lands. Det är färre
hinder till sjöss som hindrar signalen samt skillnaden i vattnets och luftens
dielektricitet gör att radiosignalen kan följa ytan lite längre innan den
klingat av.

\subsection{Rymdvåg}
\index{rymdvåg}
\index{vågutbredning!rymdvåg}

Rymdvåg är något som framför allt uppträder på kortvåg. När strålningen från
solen träffar de övre lagren i atmosfären uppstår laddade partiklar, så
kallade joner. Dessa består av fria elektroner och ibland också protoner. Fria
elektroner är ju som vi vet kännetecknet för ett ledande material så även om
luften är mycket tunn där uppe så uppstår ändå en elektriskt ledande plasma.

Ett sådant skikt kan fungera som en spegel för radiosignaler som sänds mot det
och beroende på hur kraftig joniseringen är finns det en högsta frekvens som
fortfarande reflekteras medan frekvenser över denna går igenom jonosfären och
knappt reflekteras alls. Denna högsta frekvens kallas för MUF vilket utläses
{\it Maximum Useable Frequency}, then högsta användbara frekvensen (för
jonosfärreflektion).

MUF påverkas också av ett par andra viktiga saker, dels är det dagens längd av
dygnet, sommartid kan man förvänta sig en högre MUF på grund av kraftigare
jonisering, på vintern lägre. Eftersom solen har sina solfläcksmaximum i
11-årscykler så har vi också en variation, vid solfläcksmaximum kan
joniseringen vara enormt mycket kraftigare än vid minima och detta påverkar
väldigt mycket rymdvågsutbredningen.

\subsubsection{Skipzon}
\index{vågutbrednin!skipzon}
\index{skipzon}

Beroende på frekvens och avstånd, antennhöjd, uteffekt mm så når markvågen en
viss sträcka. Rymdvågen kommer inte vara hörbar överallt utan det beror mycket
på vilken vinkel den reflekteras mot. För att komma så långt som möjligt vill
man ha en flack vinkel, kanske 10° över horisonten vilket ger en mycket
lång skipzon.

Skipzonen är den zon mellan sändaren och där rymdvågen återvänder efter att ha
reflekterats. Denna är ofta mycket längre bort än markvågen så det finns en
död zon mellan där markvågen inte orkar täcka längre och innan rymdvågen
faller in och det är denna zon som kallas för skipzonen (radiovågen ''hoppar
över denna zon'').

\subsubsection{NVIS}
\index{vågutbredning!NVIS}
\index{NVIS}

Det finns en variant av detta där man i princip eliminerar skipzonen och det
är om man nyttjar NVIS som sändningsmetod. Man placerar då antennen ganska
nära marken och på så vis riktar man större delen av energin rakt uppåt. När
denna träffar jonosfären uppstår mycket skapra vinklar som gör att signalen
återvänder över ett mycket stort område. Detta kan nyttjas och i princip
eliminera skipzonen och används ofta av militära kortvågsstationer. Med
begränsad markvåg har det också fördelen att vara mycket svårt att pejla!

NVIS är engelska och betyder \textit{Near Vertical Incidence Skywave}, alltså att man
sänder nästan vertikalt uppåt. Utbredningen kommer som regel ner i en cirkel
runt sändaren som beroende på frekvens, jonosfäriska förhållanden med mera
brukar kunna bli mellan 300--600 km.

Det betyder att man kan prata från mellansverige med större delen av
skandinavien och norra europa på NVIS.

\subsubsection{Sporadiskt E}
\index{vågutbredning!Es}
\index{sporadiskt E!Es}
\index{Es}

Ibland kan E-skiktet i jonosfären joniseras och då kan det bli riktigt fin
vågutbredning även på 6\,m bandet och förstås de övre frekvenserna på
kortvågen där man kan få riktigt bra långväga förbindelser. Sådant kallas för
''sporadiskt E'' när det uppträder och det kan vara några timmar eller över
dagen.

\subsection{Troposfärutbredning}
\index{tropo}
\index{vågutbrednin!atmosfärisk}

Under vissa väderförhållanden kan man få ett fenomen som kallas för inversion.
Normalt avtar temperaturen i atmosfären med ungefär 1$^\circ$ per 100 meter
men under dessa speciella förhållanden finner man plötsligt ett varmare
luftskikt längre upp. När detta inträffar kan radiosignalerna studsa på detta
om de kommer in med rätt vinkel eller ibland reflekteras tillbaka mot marken
på ett sätt så att markvågen kan nå mycket längre än normalt.

Detta kallas i dagligt tal för ''tropo'' och är i vissa sammanhang som för
radiolänk i tropiska länder ett stort problem men för radioamatörer är det
kära tillfällen att nå bra mycket längre än normalt på band som 6\,m, 2\,m,
och även 70\,cm.

Ibland kan signalen studsa mellan flera lager innan den kommer ut igen och
kommer ner till marken, man har då fått så kallad ''duktning'' och det går
ibland att nå mer än tiofallt den vanliga markvågsutbredningen.

Tropo uppträder främst på sommarhalvåret vid våra breddgrader.

\section{Rymdväder}
\index{Rymdväder}
\todo{Rymdväder}

Rymdvädret påverkar radiokommunikationen både i rymden och på jorden. När
solen är som mest aktiv kan det orsaka störningar på radiosystem. Detta märks
inte minst på satellitkommunikation men vid extrema fall så blir det
störningar även på markvågsförbindelser.

Vissa sorters störningar kan dock vara gynnsamma för den som använder
rymdvåg. Vid kraftigare jonisering av de olika skikten i atmosfärens övre
lager (jonosfären) så kan den börja reflektera frekvenser ända upp till
VHF-området, men mer ofta de olika kortvågsbanden, som regel 80--10 meter.

\subsection{Solens inverkan}
\index{rymdväder!solen}

\subsection{Solens cykler}
\index{solcykler}
\index{rymdväder!solcykler}

\subsection{Olika fenomen}
\index{solfenomen}
\index{rymdväder!solfenomen}

\subsubsection{Flare}
\index{flare}
\index{solstorm}
\index{rymdväder!flare}

\subsubsection{Koronamassutkastning (CME)}
\index{CME}
\index{coronal mass ejection}
\index{koronamassutkastning}
\index{rymdväder!CME}
\index{rymdväder!koronamassutkastning}

\subsubsection{Polarsken, Aurora}
\index{polarsken} \index{aurora}
\index{rymdväder!polarsken}
\index{rymdväder!aurora}

\subsection{Att förstå parametrarna}





\section{Antenner}
\index{antenner}
\todo{Antenner}

Här ska vi diskutera några av de vanligaste förekommande antennerna
och särskilt från perspektivet att få till antenner som man kan bygga
själv på egen hand till sina radioapparater. Vi kommer främst titta på
antenner för kortvåg, VHF och lite UHF (låga delen, typiskt 70 cm).


